\documentclass[10pt]{article}
\providecommand{\FIGDIR}{fig}
\newcommand{\DIGESTDATE}{11 November 2026}
\input{preamble}

\begin{document}

% ==================================================================== MASTHEAD
\thispagestyle{empty}
\begin{tikzpicture}[remember picture, overlay]
  \fill[Deep] (current page.north west) rectangle
      ($(current page.north east)+(0,-67mm)$);
  \fill[Amber] ($(current page.north west)+(0,-67mm)$) rectangle
      ($(current page.north east)+(0,-68.4mm)$);
  \foreach \x/\y/\r in {22/12/0.9, 41/26/1.5, 63/9/0.7, 88/31/1.1, 112/17/1.9,
                        138/29/0.8, 159/13/1.3, 178/54/1.0, 196/21/1.6,
                        30/44/0.6, 74/57/1.2, 125/48/0.9, 168/6/0.8, 191/59/0.7,
                        52/18/1.0, 100/60/0.6, 148/20/1.1, 66/49/1.4, 205/45/0.8}
    \fill[Sky, opacity=0.30] ($(current page.north west)+(\x mm,-\y mm)$) circle (\r mm);
\end{tikzpicture}

\vspace*{4mm}
{\sffamily\fontsize{7.6}{9}\selectfont\color{Sky}%
 AUTOMATED DAILY LITERATURE DIGEST\;\textbullet\;PhD RESEARCH WATCH\par}
\vspace{2.2mm}
{\sffamily\bfseries\fontsize{27}{29}\selectfont\color{white}%
 Particulate Matter\par}
{\sffamily\fontsize{27}{29}\selectfont\color{Sky}%
 Monitoring \& Health Effects\par}

\vspace{5.0mm}
{\sffamily\fontsize{8.4}{10}\selectfont\color{white}%
 Issue 2026-11-11\;\;\textcolor{Amber}{\textbar}\;\;Entry window 11 November 2026 (full day)\;\;%
 \textcolor{Amber}{\textbar}\;\;36 records screened in scope\par}
\vspace{18mm}

\noindent
\begin{minipage}[t]{0.190\linewidth}\metric{36}{RECORDS IN SCOPE\\(90 EXCLUDED)}{Deep}\end{minipage}\hfill
\begin{minipage}[t]{0.190\linewidth}\metric{9}{OF 10 SUBTOPIC\\BANDS POPULATED}{Teal}\end{minipage}\hfill
\begin{minipage}[t]{0.190\linewidth}\metric{6}{EFFECT ESTIMATES\\WITH A CI}{Sky}\end{minipage}\hfill
\begin{minipage}[t]{0.190\linewidth}\metric{3}{TIER-A\\RECORDS}{Amber}\end{minipage}\hfill
\begin{minipage}[t]{0.190\linewidth}\metric{10}{RECORDS CARRIED ON\\METADATA ALONE}{Clay}\end{minipage}

\vspace{6mm}

\section*{\sffamily\bfseries\color{Deep}Signal of the day}
\vspace{-1.5mm}{\color{Amber}\rule{\linewidth}{1.1pt}}\vspace{2.5mm}

\begin{keybox}
{\sffamily\bfseries\small\color{Deep}Thin the reference record to a 1-in-6 or 1-in-12 schedule and most PM$_{2.5}$ equivalence verdicts stop being determinate: across 43,576 same-site FRM--FEM pairs at 45 US sites, sparse-schedule indeterminacy reached 63.8\,\% (53.0--73.8\,\%) and 73.8\,\% (63.7--82.6\,\%), and 64.6\,\% of near-daily benchmarks were already boundary-uncertain.}\\[1.4mm]
\footnotesize
Yangtze University, \textit{Measurement Science and Technology}; 127 episodes (2022--2025), US slope--intercept equivalence envelope for the point verdict, circular moving-block bootstrap for its stability. At 1-in-3 indeterminacy was 33.0\,\% (21.1--46.1\,\%); 88 of 90 point mismatches became indeterminate rather than reversed. In a known-truth 1-in-12 simulation interval coverage was about 88\,\% and 14 of 64 ``stable'' verdicts had the wrong sign.
\textbf{Where it is weak:} one group, statistical rather than field work, and the stability rule (0.95/0.05) is a choice --- the authors show the ordering holds from 0.80 to 0.99, not the levels.
\textbf{Why it matters here:} low-cost sensor ``equivalence'' and drift-correction claims routinely rest on fewer reference days than a 1-in-6 schedule. A point verdict without its bootstrap state and reference support is not evidence of comparability.\par
\end{keybox}

\vspace{3mm}

\section*{\sffamily\bfseries\color{Deep}What changed today}
\vspace{-1.5mm}{\color{Amber}\rule{\linewidth}{1.1pt}}\vspace{2.5mm}

\footnotesize
\begin{itemize}

\item \textbf{Personal exposure, null; outdoor, a small signal.} Rural Beijing solid-fuel panel, 1,064 adults over two heating seasons with outdoor, household and 24-h personal PM$_{2.5}$ and black carbon: TNF-$\alpha$ +1.9\,\% (0.2--3.6\,\%) per 10\,$\mu$g/m$^3$ outdoor PM$_{2.5}$; nothing for household or personal exposure, or for IL-6, CRP or MDA. Household PM was imputed for $\sim$70\,\% of homes.

\item \textbf{923 kitchens on PurpleAir, no correction reported.} Eastern Ethiopia, PurpleAir II-SD for 3--5 days per home: rural 473 versus urban 80\,$\mu$g/m$^3$ (difference 393, 370--416); rural kitchens 343 (dry) to 732\,$\mu$g/m$^3$ (wet season). Uncorrected Plantower output in dense biomass smoke is humidity-biased along the very axis of the seasonal contrast.

\item \textbf{Closing the source, not measuring it.} Four California refinery retirements, difference-in-differences on births 2008--2022: Alon Bakersfield term LBW OR 0.68 (0.55--0.83), SGA 0.81 (0.74--0.89). Exposure is proximity; no pollutant was measured.

\item \textbf{Siting and geocoding.} AirQo's LOCATE ranks a context-augmented population-weighted p-median first for low-cost sensor layouts in Kampala, Jinja, Nairobi and Lagos (5--80 sensors), explicitly without any concentration field. A separate OMOP exposure toolchain shows ordinary geocoding error halves effect estimates for exposures varying over 100\,m but barely moves smooth regional ones.

\item \textbf{Exposure basis (x1).} Two records rest on low-cost sensors (the AirQo siting framework and the Ethiopian kitchens), neither producing a health estimate. Of the 13 health-side records with a recoverable basis, 62\,\% are measured, 31\,\% modelled or satellite-fused and 8\,\% proxy; of today's three effect sources only the rural Beijing TNF-$\alpha$ estimate rests on measured (village-monitor) exposure --- the glaucoma hazard ratios use a modelled surface and the refinery odds ratios a proximity proxy. \emph{No effect estimate today comes from sensor-measured exposure}.

\end{itemize}

\vspace{2mm}

\begin{keybox}
{\sffamily\bfseries\footnotesize\color{Deep}Window continuity}\\[1.2mm]
{\fontsize{7.2}{8.8}\selectfont
The 10 November issue closed the day before. This issue was built on the \textbf{12 November 23:38 UTC (17:38 CST) run}. Before 22:00 CST the newest buildable day is yesterday, so this run owes \textbf{10 and 11 November as two issues}; \textbf{11 November is covered in full}. Wednesday: no rollup is owed; \textbf{W46 (8--14 November) falls on Saturday 14 November}. A cross-day pass assigned 15 records that resurfaced today to the 10 November issue. No \textsc{backfill} records today. OpenAlex, alphaXiv, bioRxiv/medRxiv and Wiley Scholar Gateway connectors were not available on this session, and Scite was not run.\par}
\end{keybox}

\vspace{3mm}
\par\ifdim\dimexpr\pagegoal-\pagetotal\relax<18\baselineskip\newpage\fi
\section*{\sffamily\bfseries\color{Deep}Corpus analytics}
\vspace{-1.5mm}{\color{Amber}\rule{\linewidth}{1.1pt}}\vspace{2.5mm}

\noindent\begin{minipage}{\linewidth}\includegraphics[width=\linewidth]{\FIGDIR/f1_subtopics.png}
\figcap{Subtopic assignment of the 36 records across nine of ten bands: exposure assessment 13, occupational \& indoor 7, mechanistic 4, burden 3, sensing 3, other clinical 2, reproductive 2, and one each for cardiovascular/metabolic and respiratory. Neuro/mental health is empty today.}\end{minipage}

\vspace{2mm}
\noindent\begin{minipage}{\linewidth}\includegraphics[width=\linewidth]{\FIGDIR/f2_design_endpoint.png}
\figcap{Left --- architecture: 11 measurement campaigns, 10 metadata-only records, 6 modelling, 3 cohorts, 2 experimental, 2 cross-sectional, one review and one ecological (the refinery difference-in-differences). Right --- every recoverable endpoint is a singleton; the 10 metadata-only records form the only bar above one, and roughly half the endpoints are measurement or exposure outcomes rather than clinical ones.}\end{minipage}

\vspace{2mm}
\noindent\begin{minipage}{\linewidth}\includegraphics[width=\linewidth]{\FIGDIR/f3_metric_geography.png}
\figcap{Left --- particle metric as classified from the record fields: unspeciated PM or emissions 16, PM$_{2.5}$ only 12, PM$_{2.5}$ and PM$_{10}$ jointly 4, and one each for PM$_{10}$, ultrafine, size-resolved and optical properties. Right --- geography by record: China 12, North America 7, global or multi-region 7 (one review plus six metadata-only or editorial records with no stated location), Europe 4, sub-Saharan Africa 3, and one each for East Asia outside China, Oceania and South Asia.}\end{minipage}

\vspace{2mm}
\noindent\begin{minipage}{\linewidth}\includegraphics[width=\linewidth]{\FIGDIR/f6_lifecourse.png}
\figcap{Life-course windows. Preconception/in utero is the refinery-closure births and the Xiamen congenital heart defect study; childhood the metadata-only early-life obesity paper; adolescence the Seoul mobility analysis (youth subgroup, exposure only); working age the rural Beijing panel, UK Biobank, the glaucoma cohort, TiO$_2$ miners and fire instructors; older adults the Portland bedroom-sleep study (mean age $\sim$77).}\end{minipage}

\vspace{2mm}
\noindent\begin{minipage}{\linewidth}\includegraphics[width=\linewidth]{\FIGDIR/f4_heatmap.png}
\figcap{Subtopic $\times$ study architecture. Exposure assessment is the widest row, split across measurement campaigns, models and four metadata-only records; occupational \& indoor mixes measurement, a simulation and three metadata-only records. The three cohorts sit in mechanistic (rural Beijing biomarkers) and other clinical (glaucoma, UK Biobank); the single ecological cell is the reproductive refinery study.}\end{minipage}

\vspace{2mm}

\IfFileExists{\FIGDIR/extra_figs.tex}{\input{\FIGDIR/extra_figs.tex}}{}



\noindent\begin{minipage}{\linewidth}\includegraphics[width=\linewidth]{\FIGDIR/f5_forest.png}
\figcap{Ratio scale, five estimates from two studies. Three are nested follow-up horizons (30, 60, 90 days) of one glaucoma cohort, PM$_{2.5}$ increment not stated in the abstract; two are categorical post- versus pre-closure contrasts from the refinery difference-in-differences, where the exposure is proximity and no PM was measured. Different endpoints, scales of contrast and exposure definitions: no pooling and no I$^2$.}\end{minipage}

\vspace{2mm}
\noindent\begin{minipage}{\linewidth}\includegraphics[width=\linewidth]{\FIGDIR/f5_forest_2.png}
\figcap{Difference scale, one estimate: serum TNF-$\alpha$ +1.9\,\% (0.2--3.6\,\%) per 10\,$\mu$g/m$^3$ outdoor village PM$_{2.5}$, rural Beijing adults (n=1,064); the household and personal-exposure models were null and are not plotted.}\end{minipage}

\vspace{4mm}

\par\vskip 0pt plus 24\baselineskip\penalty-200\vskip 0pt plus -24\baselineskip
\section*{\sffamily\bfseries\color{Deep}Paper-level digest}
\vspace{-1.5mm}{\color{Amber}\rule{\linewidth}{1.1pt}}\vspace{2.5mm}

{\fontsize{7.4}{9}\selectfont\color{Slate}Ordered by cluster. Each entry gives the
methodological core and the finding that matters, not an abstract paraphrase. Records
marked \textit{metadata only} carry no recoverable abstract and assert no finding.\par}

\band{Cardiovascular \& metabolic\hfill 1 record}{Teal}

\paperentry{Lu Jiang Central South early-life air pollution temperature children overweight obesity evidence intervention}
{Unraveling the Effects of Early Life Air Pollution and Temperature Exposure on Children Being Overweight and Obese: Evidence and Intervention}
{Environmental Science \&amp; Technology}
{\textsc{metadata only.} Metadata only: no abstract in the Crossref deposit at issue time. Environmental Science \& Technology paper on early-life air pollution and temperature exposure and childhood overweight and obesity, with an intervention component, from Central South University. No finding is asserted here. Relevance: indirect - recheck for an abstract; potentially a tier-A early-life record once the exposure method is known.}
{10.1021/acs.est.6c13419}

\par\vskip 0pt plus 10\baselineskip\penalty-200\vskip 0pt plus -10\baselineskip
\band{Respiratory \& allergic\hfill 1 record}{Sky}

\paperentry{Christiani Thorax editorial air pollution lung cancer survival population-level benefit}
{Air pollution and lung cancer survival: moving from association to population-level benefit.}
{Thorax}
{\textsc{metadata only.} Metadata only: no abstract (editorial). Thorax editorial by Christiani on air pollution and lung cancer survival, framed around moving from association to population-level benefit; presumably accompanies a research paper in the same issue, which is not identified in the record. No finding is asserted here. Relevance: peripheral - locate the index paper in the weekly sweep.}
{10.1136/thorax-2026-225905}

\par\vskip 0pt plus 10\baselineskip\penalty-200\vskip 0pt plus -10\baselineskip
\band{Reproductive \& developmental\hfill 2 records}{Sage}

\paperentry{UCLA California four refinery retirements difference-in-differences term LBW SGA births 2008-2022}
{California petroleum refinery retirements as a natural experiment to assess effects on adverse birth outcomes.}
{Environmental epidemiology (Philadelphia, Pa.)}
{Singleton births 2008-2022 conceived within three years before or after the closure of one of four California refineries; difference-in-differences logistic regression, exposed versus unexposed communities by air district. Most retirements were followed by lower odds of term LBW and SGA; strongest for Alon Bakersfield, tLBW OR 0.68 (0.55-0.83) and SGA OR 0.81 (0.74-0.89). Exposure is proximity, not measured PM or any pollutant, so the result cannot be attributed to particles rather than VOCs, SO$_2$ or concurrent economic change; parallel trends across four distinct regions is an assumption the abstract does not test, and the per-refinery heterogeneity is not summarised. Relevance: indirect - fenceline sensor records before and after a closure would supply the exposure contrast this design lacks.}
{10.1097/ee9.0000000000000538}

\paperentry{Xiamen 890 matched pairs prenatal greenness NDVI landscape indices congenital heart defects NO$_2$ PM$_{2.5}$ mediation}
{How prenatal residential greenness links to congenital heart defects: The role of multi-pollutant mediating pathways and non-linear exposure-response.}
{Ecotoxicol Environ Saf}
{Xiamen birth records 2016-2019: 1,785 CHD cases and 1,848 controls, 890 pairs after 1:1 propensity matching. First-trimester greenness was associated with lower CHD odds; at 1 km, large core patches and connectivity were protective and fragmentation and edge density harmful. Multi-pollutant mediation found indirect paths through NO$_2$ and PM$_{2.5}$, largest for NO$_2$; the SO$_2$ path did not survive joint modelling and the O$_3$ indirect estimate was positive. The NDVI-NO$_2$ relation turned at NDVI about 0.46. No effect sizes are given in the abstract, the PM$_{2.5}$ exposure model is not stated, and propensity matching discarded half the cases. Relevance: indirect - PM$_{2.5}$ enters only as a mediator of greenness.}
{10.1016/j.ecoenv.2026.121023}

\par\vskip 0pt plus 10\baselineskip\penalty-200\vskip 0pt plus -10\baselineskip
\band{Mechanistic toxicology\hfill 4 records}{Violet}

\paperentry{Central China Normal mouse 12-week PM$_{2.5}$ liver lipid accumulation IRE1a-XBP1s ER stress inhibitor}
{PM$_{2.5}$-induced liver injury mediated by lipid accumulation via the endoplasmic reticulum stress-related IRE1alpha-XBP1s pathway in mice.}
{J Environ Sci (China)}
{Mice exposed to PM$_{2.5}$ for 12 weeks; histopathology, liver-injury markers and proteomics. Lipid accumulation stemmed mainly from triglyceride synthesis; differentially expressed liver proteins clustered in ER processing and fatty-acid metabolism, and IRE1a-XBP1s pathway protein expression rose with PM$_{2.5}$ concentration. Co-administration of a specific IRE1a-XBP1s inhibitor lowered pathway expression and attenuated both lipid accumulation and liver injury. The abstract gives no exposure concentration, route, particle source or animal numbers, so the dose cannot be placed against human ambient levels, and the inhibitor result supports necessity of the pathway in this model only. Relevance: indirect - adds a hepatic steatosis mechanism to the metabolic arm of PM toxicity; no measurement content.}
{10.1016/j.jes.2026.02.006}

\paperentry{Rural Beijing solid-fuel cohort 1064 adults outdoor household personal PM$_{2.5}$ black carbon TNF-alpha IL-6 CRP MDA}
{Exposure-response relationships of outdoor and household air pollution with biomarkers of inflammation and oxidative stress in rural Chinese adults.}
{Environ Res}
{Longitudinal panel of 1,064 adults (mean age 61) in 50 villages over two heating seasons (2018-19, 2019-20), in a region where residential solid fuel dominates PM$_{2.5}$. Outdoor PM$_{2.5}$ was monitored continuously per village, household PM$_{2.5}$ in about 30\% of homes, plus 24-h personal PM$_{2.5}$ and black carbon; models with individual random effects. TNF-alpha rose 1.9\% (95\% CI 0.2-3.6\%) per 10 $\mu$g\,m$^{-3}$ outdoor PM$_{2.5}$; no clear association for household or personal exposure, or with IL-6, CRP or MDA. Household exposure was multiply imputed for \textasciitilde{}70\% of homes, which biases those estimates toward the null, and four biomarkers by four exposures invites one chance hit. Relevance: direct - one of few cohorts with outdoor, indoor and personal PM side by side; the null for measured personal exposure is the result to watch.}
{10.1016/j.envres.2026.126092}

\paperentry{Luxembourg Institute of Health ExposomeMap-PM disease map 304 entities 497 processes AOP transcriptomic overlay}
{ExposomeMap-PM: fine and ultrafine particle-induced toxicity - from adverse outcome pathways to systems biomedicine disease models.}
{Front Toxicol}
{A curated mechanistic disease map linking fine and ultrafine particle exposure to outcomes: 304 entities and 497 processes in five blocks (alveolar deposition, oxidative stress, barrier disruption, mitochondrial dysfunction, pro-inflammatory signalling) ending in asthma, COPD, atherosclerosis and neurodegeneration, built in SBGN and served on MINERVA. In a test overlay from one reference particle and a 3-day exposure, map nodes captured 1.50-fold more differentially expressed transcripts than size-matched random gene sets (permutation p=0.0074); proteomics 1.13-fold. One particle type and a short exposure validate coverage, not completeness, and enrichment is partly guaranteed by curating from the same literature. Relevance: indirect - a reusable scaffold for attaching composition-resolved exposure, which low-cost mass sensors do not provide.}
{10.3389/ftox.2026.1951777}

\paperentry{Laval retinal pigment epithelium wood smoke extract high-energy visible blue light synergy Nrf2 oxidative stress}
{Toxic synergy between wood combustion smoke components and high energy visible light on retinal pigment epithelial cells.}
{J Photochem Photobiol B}
{Human retinal pigment epithelium cells co-exposed to wood combustion smoke extract and high-energy visible (blue) light. The two acted synergistically to reduce viability through oxidative stress that antioxidant treatment mitigated, implying light-sensitive chromophores among smoke constituents; cells activated Nrf2 detoxification. No doses, extract preparation, or how smoke components would reach the retina in vivo are given in the abstract, and an extract is not deposited particulate matter. The plausible exposure path is ocular surface contact or systemic translocation, neither of which the design addresses; the result is a hazard signal for heavy wood-smoke settings rather than an ambient effect. Relevance: peripheral - extends wood-smoke toxicity to a photosensitisation route; no exposure-measurement content.}
{10.1016/j.jphotobiol.2026.113461}

\par\vskip 0pt plus 10\baselineskip\penalty-200\vskip 0pt plus -10\baselineskip
\band{Exposure assessment \& modelling\hfill 13 records}{Clay}

\paperentry{Urumqi-Changji-Shihezi soil moisture PM$_{2.5}$ PM$_{10}$ sensitivity threshold CALIOP aerosol typing 3D}
{Three-dimensional identification of dominant atmospheric pollutants and their sensitivity to soil moisture in northwestern China urban agglomerations.}
{J Environ Sci (China)}
{Northwest China urban agglomeration: surface PM$_{2.5}$ and PM$_{10}$ related to soil moisture, with CALIOP lidar used to map aerosol subtypes in three dimensions. Sensitivity thresholds of both fractions to soil moisture fell at 0.2266-0.2389 m3/m3 (precipitation effects excluded), PM$_{10}$ being more sensitive. Moist mountain units showed negative PM-moisture correlation (suppressed dust suspension), dry desert and plain units positive correlation (resuspension). Polluted dust dominated below 4 km, polluted continental aerosol below 2 km, with a stable mixed-dust layer aloft. The abstract does not state the PM data source, period or how thresholds were estimated. Relevance: indirect - dust-dominated regimes are where optical low-cost sensors under-read coarse mass; a moisture threshold is a usable covariate for that correction.}
{10.1016/j.jes.2026.04.026}

\paperentry{Steel plant CCSEM 91185 single particles Na2SO4 sulfate near-source sedimentation transport 6-20 km}
{Revealing the contribution of sulfate-containing particles from industrial emissions to atmospheric environment through CCSEM: Insight into morphology, formation, and transport.}
{J Environ Sci (China)}
{CCSEM classified 91,185 individual particles from five sites within a steel plant and its surroundings. By number, Na-rich (46.6\%), Ca-rich (21.5\%) and Si-rich (8.2\%) particles dominated, spanning 0.2-43.6 um. Elongated submicron Na2SO4 particles made up 72.8\% of particles nearest the exhaust outlet, aged rapidly (mean size 1.1 to 3.4 um) with 96.2\% depositional removal inside the plant, and fell to 5.3\% (2.2-11.7\%) at ambient sites 6-20 km away. Number fractions are not mass fractions, a single plant and campaign are shown, and the deposition figure is inferred from composition gradients rather than flux measurement. Relevance: indirect - fenceline sensor networks around steel plants will see a size distribution very different from the urban background they were calibrated in.}
{10.1016/j.jes.2026.03.096}

\paperentry{Seoul 424 dongs hourly PM$_{2.5}$ surface mobile-network population static vs dynamic exposure discrepancy 2023}
{Hourly population movement and static-to-dynamic PM$_{2.5}$ exposure discrepancies in Seoul, South Korea.}
{Environ Pollut}
{All 424 Seoul administrative neighbourhoods in 2023: residence-based daily-mean PM$_{2.5}$ compared with a place-based estimator weighting hourly concentrations by mobile-network population counts. The citywide mean daily discrepancy was modest but systematic: residence weighting underestimated working-hour exposure in activity hubs and evening/overnight exposure in residential neighbourhoods, while mixed-use areas sat near zero. Discrepancies were larger on weekdays and in winter and spring; working-age adults had the largest gaps in hubs, but winter gaps among youth were 25-51\% greater. Seasonal variation reflected PM$_{2.5}$ levels, not population rhythm. Both inputs are modelled surfaces, and indoor time is ignored. Relevance: direct - hourly sensor networks are what make the place-based estimator possible; daily-mean surfaces discard exactly this signal.}
{10.1016/j.envpol.2026.129480}

\paperentry{Sichuan Basin December 2018 dust intrusion WRF-Chem aerosol-radiation interaction cross-mountain PM$_{10}$ ozone}
{Aerosol-radiation interactions reshape cross-mountain dust transport and downwind air quality responses over a mountain-basin system.}
{J Hazard Mater}
{One dust intrusion into the Sichuan Basin (2-5 December 2018) analysed with observations, Lagrangian dispersion and WRF-Chem with and without aerosol-radiation interactions. Dust crossed the Qinling-Daba mountains on elevated pathways, raising PM$_{10}$ and ozone. ARI cut maximum boundary-layer height by 17.5\% over the plateau and 28.7\% in the basin, lowered cumulative cross-mountain PM$_{10}$ transport by 13.6\% and near-surface basin PM$_{10}$ by 23.5 $\mu$g\,m$^{-3}$ (about 10\%); ozone rose 8.2 ppbv. A single event and a single model configuration; the ARI effect depends on the dust optical properties assumed. Relevance: indirect - forecasts that omit ARI overpredict basin PM$_{10}$ during dust events, which matters when they are used to flag sensor outliers.}
{10.1016/j.jhazmat.2026.144174}

\paperentry{Brown Rhode Island 144 sites 531 site-sessions PM$_{2.5}$ sound night-time illuminance variance place vs time}
{The acoustic environment is structured more by place and fine particulate matter more by time across Rhode Island communities}
{Research Square}
{Preprint. PM$_{2.5}$, sound and night-time illuminance measured at 144 sites in eight Rhode Island municipalities, four day/night session types over 39 days (531 site-sessions). Sound and PM$_{2.5}$ were uncorrelated and had different predictors: sound rose near industrial land and varied mainly between sites; PM$_{2.5}$ tracked weather (falling with wind, rising with temperature and precipitation) and varied mainly between days, rising at night. Results held under penalised regression and spatial-cluster adjustment. The instrument is not named, and short sessions at one site per visit confound site with day, which is exactly the decomposition being estimated. Relevance: direct - quantifies how little between-site PM$_{2.5}$ contrast short mobile campaigns resolve against day-to-day variance.}
{10.21203/rs.3.rs-10385808/v2}

\paperentry{Ojota landfill Lagos bio-cover QFIELD QGIS PM$_{2.5}$ PM$_{10}$ SO$_2$ CO six ML methods AQI risk surface}
{Machine Learning-Enhanced Air Quality Management via QFIELD and QGIS: Mapping and Bio-cover Materials at Ojota Landfill}
{Research Square}
{Preprint. Six geo-referenced sampling points at 15 m spacing across the Ojota landfill, Lagos, measuring PM$_{2.5}$, PM$_{10}$, SO$_2$ and CO as baseline, then six machine-learning methods (random forest, gradient boosting, SVM, kNN, Gaussian process/regression kriging, ANN) to classify AQI drivers and produce an uncertainty-aware risk surface; a laboratory bio-cover (topsoil, compost, rice husk, earthworm castings) reduced microbial indicators. No PM concentrations, instrument or model skill are reported in the abstract, and six points cannot support six learned spatial models. Relevance: peripheral - shows the demand for landfill-fenceline PM monitoring in West Africa; the analysis is underpowered.}
{10.21203/rs.3.rs-11369151/v1}

\paperentry{Tufts OHDSI GIS toolchain OMOP exposure infrastructure synthetic 40000 persons simulated PM$_{2.5}$ geocoding error attenuation}
{An OMOP-Compatible Exposure Infrastructure for GeoAI-Enabled Environmental Health Research: The OHDSI GIS Toolchain}
{Research Square}
{Preprint describing an open-source toolchain that standardises geospatial sources into person-place-time exposure records in the OMOP Common Data Model (FAIR catalogue, GIS vocabulary, local geocoding, geometry-attribute index). Tested on a synthetic cohort of 40,000 persons placed in real US census tracts with eight simulated PM$_{2.5}$ effects and six nulls: derived exposures matched an independent answer key and known effects were recovered. Displacing residences by empirical geocoding error changed estimates little for smooth regional exposure but attenuated them by up to half for exposures varying over 100 m. Synthetic data test plumbing, not real-world exposure validity. Relevance: direct - quantifies why fine-scale (sensor-resolved) surfaces lose half their signal under ordinary geocoding error.}
{10.21203/rs.3.rs-11436369/v1}

\paperentry{Wadden Sea size-resolved inorganic ions levoglucosan 13C 14C organic carbon ship sulfate continental outflow}
{Measurement report: Size-resolved seasonal study of inorganic ions and isotopic carbon signatures of aerosol particles at the Wadden Sea}
{Atmospheric Chemistry and Physics}
{Measurement report from a Wadden Sea coastal site: size-resolved samples by season and wind sector analysed for inorganic ions, levoglucosan, total carbon and stable and radiocarbon isotopes of organic carbon at several desorption temperatures. Land-sector samples were dominated by nitrate, ammonium and total carbon, the single sea-sector set by sulfate and total carbon; sulfate reached 0.5 $\mu$g\,m$^{-3}$ in both, suggesting ship emissions. Fossil OC fractions were highest below 250 nm (traffic combustion), and 13C signatures pointed to SOA in smaller, more volatile carbon and photochemical processing at larger sizes. One sea-sector sample set limits the ship attribution. Relevance: peripheral - background composition for coastal network sites; no sensing content.}
{10.5194/acp-26-15931-2026}

\paperentry{Cholakian Bertrand northern hemisphere multi-annual concentration exposure simulations analysis}
{Analysis of large-scale multi-annual concentration and exposure simulations for the northern hemisphere}
{Atmospheric Environment}
{\textsc{metadata only.} Metadata only: no abstract in the Crossref deposit at issue time. Atmospheric Environment paper analysing large-scale, multi-annual concentration and exposure simulations over the northern hemisphere (first author Cholakian). No finding is asserted here. Relevance: indirect - hemispheric model exposure fields are the background that regional sensor fusion corrects; recheck in the weekly sweep.}
{10.1016/j.atmosenv.2026.122445}

\paperentry{Jana Delhi road waterlogging post-event PM$_{2.5}$ satellite-based association robustness road network}
{Road waterlogging and post-event PM$_{2.5}$ on Delhi's road network: a satellite-based assessment of association and its robustness}
{Atmospheric Pollution Research}
{\textsc{metadata only.} Metadata only: no abstract in the Crossref deposit at issue time. Atmospheric Pollution Research paper using satellite data to test an association between road waterlogging and post-event PM$_{2.5}$ on Delhi's road network, with robustness checks. No finding is asserted here. Relevance: indirect - a resuspension mechanism that kerbside low-cost networks in Delhi would observe directly; recheck for an abstract.}
{10.1016/j.apr.2026.103270}

\paperentry{Zhang Tang Wuhan PM$_{2.5}$ temporal variation water-soluble ions polar organic compounds driving factors}
{Temporal variations and driving factors of PM$_{2.5}$ pollution in Wuhan, central China: Insights from water-soluble ions and polar organic compounds}
{Atmospheric Pollution Research}
{\textsc{metadata only.} Metadata only: no abstract in the Crossref deposit at issue time. Atmospheric Pollution Research paper on temporal variation and driving factors of PM$_{2.5}$ in Wuhan from water-soluble ions and polar organic compounds. No finding is asserted here. Relevance: peripheral - composition and source context for central China; recheck for an abstract.}
{10.1016/j.apr.2026.103269}

\paperentry{Tibetan Plateau QOMS NAMORS particle-bound organophosphate esters PM$_{2.5}$ warm-season maximum trans-Himalayan transport}
{Contrasting Dynamics of Particle-Bound Organophosphate Esters Reveal Distinct High-Altitude Contaminant Transport Processes over the Tibetan Plateau}
{Environmental Science \&amp; Technology}
{Year-round synchronous PM$_{2.5}$ sampling at a southern Himalayan receptor (QOMS) and an inland background site (NAMORS). Sum of OPEs averaged 111 (SD 125) pg/m3 at QOMS and 20 (17) at NAMORS, with a counterintuitive warm-season maximum coinciding with a 72-91\% fall in hygroscopic secondary inorganic aerosol, read as trans-Himalayan monsoon transport persisting despite wet deposition. Chlorinated congeners (TCPP, TDCPP) dominated (TCEP/TCPP 0.09-0.46), unlike Arctic profiles; trajectories and PCA point to South Asian sources. The authors note the absence of simultaneous gas-phase data, which leaves partitioning unconstrained. Relevance: peripheral - toxic trace constituents invisible to any mass sensor; remote background context.}
{10.1021/acs.est.6c11837}

\paperentry{Tang Zhao spatiotemporal varying-coefficient calibration heterogeneous measurement error gridded PM$_{2.5}$ exposure estimates}
{Accounting for Spatiotemporally Heterogeneous Measurement Error in Gridded PM$_{2.5}$ Exposure Estimates: A Spatiotemporal Varying-Coefficient Calibration Approach}
{Environmental Pollution}
{\textsc{metadata only.} Metadata only: no abstract in the Crossref deposit at issue time. Environmental Pollution paper proposing a spatiotemporal varying-coefficient calibration to account for spatiotemporally heterogeneous measurement error in gridded PM$_{2.5}$ exposure estimates. No finding is asserted here. Relevance: direct - error-aware calibration of exposure surfaces against monitors is the step where low-cost network data enter; priority recheck.}
{10.1016/j.envpol.2026.129538}

\par\vskip 0pt plus 10\baselineskip\penalty-200\vskip 0pt plus -10\baselineskip
\band{Sensing, forecasting \& instrumentation\hfill 3 records}{Amber}

\paperentry{AirQo LOCATE GeoAI population-weighted p-median sensor placement Kampala Jinja Nairobi Lagos budgets 5-80}
{LOCATE: Context-Aware GeoAI Decision Support for Population-Oriented Air-Quality Sensor Placement in Data-Constrained Cities}
{Research Square}
{Preprint from the AirQo group. LOCATE is a pre-deployment siting framework for data-constrained cities that prioritises proximity to modelled residential population while adding road and building context. Six feature-informed clustering methods, two allocation methods and two uninformed baselines were compared across Kampala, Jinja, Nairobi and Lagos for budgets of 5-80 sensors (20 seeds). The context-augmented population-weighted p-median ranked first in all four cities; coverage rose with budget but stayed lower in the larger Lagos and Nairobi domains. The authors state the limitation plainly: layouts are shortlists, not evidence of pollution reconstruction or exposure accuracy, because no concentration field enters the objective. Relevance: direct - siting for low-cost PM networks where no reference surface exists.}
{10.21203/rs.3.rs-11439350/v1}

\paperentry{IIT Madras operating procedures advanced aerosol instruments high-altitude observatories technical note}
{Effective operating procedures for advanced aerosol instruments at high-altitude observatories: A technical note}
{Aerosol Science and Technology}
{\textsc{metadata only.} Metadata only: no abstract in the Crossref deposit at issue time. A technical note in Aerosol Science and Technology on operating procedures for advanced aerosol instruments at high-altitude observatories, from IIT Madras. No finding is asserted here. Relevance: indirect - high-altitude sites are the reference anchors for regional aerosol networks; recheck for an abstract in the weekly sweep.}
{10.1080/02786826.2026.2755497}

\paperentry{Yangtze 45 US sites 43576 reference-FEM PM$_{2.5}$ pairs sparse 1-in-3 1-in-6 1-in-12 sampling comparability verdict bootstrap indeterminacy}
{Decision uncertainty in particulate-matter measurement comparability under sparse reference sampling}
{Measurement Science and Technology}
{45 US sites (2022-2025), 127 episodes, 43,576 same-site reference-equivalent PM$_{2.5}$ pairs. Each near-daily record was thinned to every phase of 1-in-3, 1-in-6 and 1-in-12 schedules; US equivalence slope-intercept envelopes gave point verdicts, a moving-block bootstrap their stability. 64.6\% of near-daily benchmarks were already boundary-uncertain; on the 45 stable episodes, sparse-schedule indeterminacy was 33.0\% (21.1-46.1\%), 63.8\% (53.0-73.8\%) and 73.8\% (63.7-82.6\%), and 88 of 90 point mismatches became indeterminate rather than reversed. In a known-truth 1-in-12 simulation, coverage was about 88\% and 14 of 64 stable verdicts had the wrong sign. Single-author-group statistical work, not yet replicated. Relevance: direct - most low-cost sensor 'equivalence' claims rest on sparser reference support than this; report bootstrap state with the verdict.}
{10.1088/1361-6501/aec0bc}

\par\vskip 0pt plus 10\baselineskip\penalty-200\vskip 0pt plus -10\baselineskip
\band{Occupational \& indoor\hfill 7 records}{Amber}

\paperentry{Western Australia mining TiO2 604 personal inhalable samples 2018-2025 proposed ultrafine fine exposure limits}
{Titanium dioxide exposure in the Western Australian mining industry: Implications for compliance with current and proposed exposure limits.}
{J Occup Environ Hyg}
{De-identified regulator data: 604 personal inhalable TiO2 samples from Western Australian mining, 2018-2025, compared with the current 10 mg/m3 inhalable limit and proposed limits of 0.3 mg/m3 (ultrafine/nanoscale) and 2.4 mg/m3 (fine, respirable). Exposures were consistently low; one sample exceeded the numerical value of the proposed ultrafine limit and the second highest was 0.12 mg/m3; none exceeded 2.4 mg/m3; trends were stable or falling. The weakness is the conclusion itself: inhalable gravimetric mass says nothing about the ultrafine/nanoscale fraction the 0.3 mg/m3 limit targets, so 'conservative evidence of compliance' does not follow for that limit. Relevance: peripheral - illustrates why size-specific limits need size-resolved measurement, not legacy inhalable samples.}
{10.1080/15459624.2026.2749875}

\paperentry{Haramaya HDSS 923 kitchens PurpleAir II-SD PM$_{2.5}$ rural biomass 473 vs urban 80 $\mu$g\,m$^{-3}$ wet-dry season}
{PM$_{2.5}$ Concentrations among Rural Biomass-Fuel and Urban Mixed-Fuel Households in Eastern Ethiopia: A Comparative Cross-Sectional Study}
{Research Square}
{Preprint. 923 kitchens (573 rural, 350 urban) in the Haramaya HDSS, eastern Ethiopia, monitored for 3-5 days each with 20 PurpleAir II-SD sensors, January 2023 to December 2024. Mean 24-h PM$_{2.5}$ was 473 (SD 255) $\mu$g\,m$^{-3}$ in rural and 80 (92) in urban kitchens, difference 393 $\mu$g\,m$^{-3}$ (95\% CI 370-416). In 96 kitchens measured in both seasons, rural means rose from 343 (dry) to 732 $\mu$g\,m$^{-3}$ (wet), urban from 47 to 105. All rural homes used open fires. The abstract reports no correction or co-location: uncorrected Plantower readings in dense biomass smoke at high humidity carry large, humidity-dependent bias, which is the same axis as the wet-season contrast. Relevance: direct - the largest low-cost household PM dataset in the Horn of Africa; usable only once a smoke- and RH-specific correction is shown.}
{10.21203/rs.3.rs-11347383/v1}

\paperentry{Thompson Jung airborne microplastic generation transport 3D printing university setting}
{Airborne Microplastic Generation and Transport from 3D Printing Activities in a University Setting}
{Aerosol and Air Quality Research}
{\textsc{metadata only.} Metadata only: no abstract in the Crossref deposit at issue time. Aerosol and Air Quality Research paper on airborne microplastic generation and transport from 3D-printing activities in a university setting. No finding is asserted here. Relevance: indirect - 3D printers are a known indoor ultrafine source that consumer optical sensors under-count; recheck for an abstract.}
{10.1007/s44408-026-00180-9}

\paperentry{Xiangtan cleanroom hand activity particle deposition model airflow velocity particle size height}
{Study on the Influence of Personnel Hand Activities in a Cleanroom on the Deposition Characteristics of Released Particulate Matter}
{Environmental Pollution}
{A particle deposition model for operator hand movements in cleanrooms, checked against existing experimental data and run numerically over hand height, particle size and supply-air velocity. Deposition of 5 and 10 um particles far exceeded that of 0.5-2 um; size differences narrowed with airflow, the 10:5 um deposition-fraction ratio falling from 3.2 at 0.25 m/s to 1.25 at 0.45 m/s; deposition fell sharply up to 0.15 m hand height and slowly above. Inhalation is not modelled despite the occupational framing, and validation data are not described. Relevance: peripheral - near-source coarse-particle deposition; no sensing content.}
{10.1016/j.envpol.2026.129490}

\paperentry{Vicente Rysavy bioethanol fireplace fuels burners particulate-bound organics black carbon morphology toxicity}
{Particulate-bound organic compounds, black carbon, morphology and toxicity from a bioethanol fireplace operated with different fuels and burners}
{Building and Environment}
{\textsc{metadata only.} Metadata only: no abstract in the Crossref deposit at issue time. Building and Environment paper characterising particulate-bound organic compounds, black carbon, particle morphology and toxicity from a bioethanol fireplace across fuels and burners. No finding is asserted here. Relevance: indirect - unflued indoor combustion is a growing indoor PM source; recheck for an abstract.}
{10.1016/j.buildenv.2026.115577}

\paperentry{UL FSRI NIOSH live-fire training structures exterior real-time optical PM instructors fuels location}
{Potential Airborne Particulate Matter (PM) Exposures for Firefighters Working Outside Live-Fire Training Structures.}
{Fire Technol}
{Stationary optical real-time PM area monitors outside two training structures: a single-storey fire-behaviour lab (five fuels) and a two-storey concrete multicompartment structure (three fuels). In the fire-behaviour lab, fuel load had limited effect on exterior PM, but standing upwind and further from the structure lowered the hazard; at the multicompartment structure fuel load mattered more and the second floor showed high PM where instructors may be unable to relocate. No concentrations, instrument model or smoke-specific correction are given in the abstract; area monitors are not personal breathing-zone exposure. Relevance: indirect - optical monitors in fresh smoke need a smoke-specific mass correction before such numbers are compared with limits.}
{10.1007/s10694-026-02032-9}

\paperentry{OHSU Portland 29 older adults bedroom PM$_{2.5}$ light noise temperature humidity TVOC CO$_2$ objective sleep 211 nights}
{Associations between night-time bedroom environmental factors and objective sleep metrics of older adults.}
{Int Conf Indoor Air Qual Vent Energy Conserv Build}
{29 older adults (17 cognitively normal, 12 with MCI or dementia; mean age about 77) followed for a mean 211 nights at home with objective sleep metrics and seven bedroom environmental measures. Light had the strongest association with total sleep time, humidity with WASO and efficiency, temperature with movement and heart and respiratory rate; PM$_{2.5}$ was associated with total sleep time (effect size -0.050) and sleep efficiency (-0.058), Bonferroni significant. The PM$_{2.5}$ sensor is not identified, effect sizes are standardised without stated units, and 29 people give many nights but little between-person information. Conference paper. Relevance: direct - long-run in-home PM sensing paired with passive physiology is the deployment model for indoor exposure-response work.}
{10.1051/e3sconf/202671605024}

\par\vskip 0pt plus 10\baselineskip\penalty-200\vskip 0pt plus -10\baselineskip
\band{Burden, policy \& mitigation\hfill 3 records}{Coral}

\paperentry{Hangzhou Asian Games emission controls WSOC FT-ICR MS organosulfate oligomers quasi-natural experiment 2026}
{Significant suppression of sulfur-containing organic aerosols by emission controls during the Hangzhou Asian Games.}
{J Environ Sci (China)}
{Quasi-natural experiment around the 19th Asian Games: WSOC in PM$_{2.5}$ characterised by ultrahigh-resolution FT-ICR MS across control and post-Games periods. Controls coincided with selective depletion of high-molecular-weight organosulfate oligomers (O10-O13) and highly oxygenated CHOS species, consistent with suppressed acid-catalysed oligomerisation and curtailed ageing; Games-period pollution was dominated by humidity-enhanced aqueous processing, the post-Games rebound by gas-phase oxidation yielding organic dinitrates. No PM$_{2.5}$ mass change, sample count or uncertainty is given in the abstract, and a single event with co-varying meteorology cannot separate control effects from weather. Relevance: indirect - composition shifts of this kind change optical-sensor response per unit mass, a reason not to transfer calibrations across policy regimes.}
{10.1016/j.jes.2026.03.042}

\paperentry{Barcelona Granada Palma Valencia 100 m urban proximity index PM$_{2.5}$ NO$_2$ NDVI HIA attributable mortality quartiles}
{The interplay of urban proximity, environmental exposures and health impacts in Spanish cities.}
{Environ Pollut}
{Four Spanish cities: a 100 m urban proximity index (walkability, cyclability, vitality, 15-minute city) overlaid on PM$_{2.5}$, NO$_2$, NDVI and population, with grid-level health impact assessment of mortality from non-compliance with WHO guidelines. The most proximate quartile carried the highest burden: PM$_{2.5}$-attributable mortality 34-102 versus 21-87 per 100,000 in the least proximate, NO$_2$ 58-154 versus 20-121, low greenness 21-35 versus 8-17. Ranges span cities; no uncertainty intervals are reported, and applying one concentration-response function across quartiles assumes away the within-city confounding the proximity index carries. Relevance: indirect - central, compact districts are where a 100 m exposure surface most needs dense ground observation to be credible.}
{10.1016/j.envpol.2026.129400}

\paperentry{Chatoutsidou Lazaridis European major cities mortality burden long-term PM$_{2.5}$ NO$_2$ updated evidence}
{Mortality burdens from long-term exposure to PM$_{2.5}$ and NO$_2$ in major European cities: what does updated evidence suggest?}
{Atmospheric Pollution Research}
{\textsc{metadata only.} Metadata only: no abstract in the Crossref deposit at issue time. Atmospheric Pollution Research paper on mortality burdens from long-term PM$_{2.5}$ and NO$_2$ exposure in major European cities under updated concentration-response evidence (first author Chatoutsidou). No finding is asserted here. Relevance: indirect - burden estimates of this kind inherit the exposure surface; recheck for an abstract.}
{10.1016/j.apr.2026.103272}

\par\vskip 0pt plus 10\baselineskip\penalty-200\vskip 0pt plus -10\baselineskip
\band{Other clinical endpoints\hfill 2 records}{Slate}

\paperentry{Fudan three-centre 7820 eyes incisional glaucoma surgery early failure PM$_{2.5}$ PM$_{10}$ NO$_2$ O$_3$ CO Cox 2015-2025}
{Association of Ambient Air Pollution Exposure With Early Failure After Incisional Glaucoma Surgery.}
{Ophthalmol Glaucoma}
{Three tertiary centres, 7,820 eyes after incisional glaucoma surgery (2015-2025), 466 failures by 30 days, 711 by 60 and 1,083 by 90. PM$_{2.5}$ was associated with failure at every horizon: adjusted HR 1.090 (1.001-1.187) at 30 days, 1.141 (1.064-1.224) at 60 and 1.240 (1.162-1.324) at 90, roughly linear on splines; O$_3$ and CO mattered early, PM$_{10}$ and NO$_2$ later, SO$_2$ not at all; no between-centre heterogeneity. The abstract does not state the exposure increment or window, and six pollutants across three horizons with subgroup interactions invite multiplicity. Surgical failure has strong clinical determinants (prior surgery, antimetabolite use) only partly captured. Relevance: peripheral - an unusual endpoint; the exposure is a modelled residential surface.}
{10.1016/j.ogla.2026.11.001}

\paperentry{UK Biobank 348247 PM$_{2.5}$ PM$_{10}$ PM$_{2.5}$-10 absorbance NO$_2$ oral inflammatory disease KDM-BA PhenoAge mediation}
{Associations between air pollution, accelerated biological aging, and oral inflammatory diseases: evidence from the UK Biobank cohort.}
{Front Public Health}
{UK Biobank, 348,247 adults, median follow-up 16.49 years, 3,967 incident oral inflammatory diseases. Per SD increment, PM$_{2.5}$, PM$_{10}$, PM$_{2.5}$-10 and PM$_{2.5}$ absorbance were associated with 8\%, 4\%, 3\% and 5\% higher incidence (NO$_2$ 10\%, NOx 8\%; all p<0.05). PhenoAge acceleration mediated 8.52-16.65\% of the associations, KDM biological age 1.21-5.42\%. The abstract reports no confidence intervals. Exposure is a single-year (2010) land-use regression surface applied over 16 years, the outcome comes from hospital coding of a condition mostly managed in primary dental care, and mediation by a biomarker measured at baseline alongside exposure cannot establish temporal ordering. Relevance: peripheral - modelled exposure, no measurement content.}
{10.3389/fpubh.2026.1869182}

\clearpage
\section*{\sffamily\bfseries\color{Deep}Corpus register}
\vspace{-1.5mm}{\color{Amber}\rule{\linewidth}{1.1pt}}\vspace{2.5mm}

\input{table_rows}

\vspace{2mm}

\begin{keybox}
{\sffamily\bfseries\footnotesize\color{Deep}Provenance and method}\\[1.2mm]
{\fontsize{7.2}{8.8}\selectfont
Entry window \textbf{11 November 2026}, single day (Wednesday). Harvester legs: \textbf{PubMed} health \textbf{31}, sensing \textbf{3}; \textbf{Europe PMC} \textbf{41}; \textbf{Crossref by ISSN} \textbf{71}; \textbf{arXiv} no in-window entry; \textbf{OpenAlex} HTTP 429 on the harvester leg. 146 raw, 138 unique, 5 already carried, 133 fresh. PubMed connector (EDAT 10--11 November, three axes) 60 PMIDs (\texttt{\_conn\_1112.py}), 31 bucketed to 11 November, 0 seen, 2 PMIDs merged onto existing records, 8 appended (141); cross-day dedupe moved 15 records to the 10 November issue --- \textbf{126 screened}. Abstracts: efetch filled 35 of 36 PMID records; Crossref/Semantic Scholar recovered none of ten Crossref records tried.
OpenAlex, alphaXiv, bioRxiv/medRxiv and Wiley Scholar Gateway connectors were not available on this session; \textbf{Scite} retraction check not run, so no record carries a Scite check. No Amass or Consensus backfill this run.
\textbf{36 admitted}, \textbf{90 rejected} with logged reasons. \textbf{10 carried on metadata alone}. Affiliations from PubMed XML and Crossref (OpenAlex unavailable); blank where neither carries one. Trial watch: see W46 weekly.
Abstracts remain the copyright of their publishers.\par}
\end{keybox}

\end{document}
